% Options for packages loaded elsewhere
\PassOptionsToPackage{unicode}{hyperref}
\PassOptionsToPackage{hyphens}{url}
\PassOptionsToPackage{dvipsnames,svgnames,x11names}{xcolor}
\documentclass[
  11pt,
]{article}
\usepackage{xcolor}
\usepackage[margin=25mm]{geometry}
\usepackage{amsmath,amssymb}
\setcounter{secnumdepth}{-\maxdimen} % remove section numbering
\usepackage{iftex}
\ifPDFTeX
  \usepackage[T1]{fontenc}
  \usepackage[utf8]{inputenc}
  \usepackage{textcomp} % provide euro and other symbols
\else % if luatex or xetex
  \usepackage{unicode-math} % this also loads fontspec
  \defaultfontfeatures{Scale=MatchLowercase}
  \defaultfontfeatures[\rmfamily]{Ligatures=TeX,Scale=1}
\fi
\usepackage{lmodern}
\ifPDFTeX\else
  % xetex/luatex font selection
\fi
% Use upquote if available, for straight quotes in verbatim environments
\IfFileExists{upquote.sty}{\usepackage{upquote}}{}
\IfFileExists{microtype.sty}{% use microtype if available
  \usepackage[]{microtype}
  \UseMicrotypeSet[protrusion]{basicmath} % disable protrusion for tt fonts
}{}
\makeatletter
\@ifundefined{KOMAClassName}{% if non-KOMA class
  \IfFileExists{parskip.sty}{%
    \usepackage{parskip}
  }{% else
    \setlength{\parindent}{0pt}
    \setlength{\parskip}{6pt plus 2pt minus 1pt}}
}{% if KOMA class
  \KOMAoptions{parskip=half}}
\makeatother
\usepackage{listings}
\newcommand{\passthrough}[1]{#1}
\lstset{defaultdialect=[5.3]Lua}
\lstset{defaultdialect=[x86masm]Assembler}
\usepackage{longtable,booktabs,array}
\usepackage{calc} % for calculating minipage widths
% Correct order of tables after \paragraph or \subparagraph
\usepackage{etoolbox}
\makeatletter
\patchcmd\longtable{\par}{\if@noskipsec\mbox{}\fi\par}{}{}
\makeatother
% Allow footnotes in longtable head/foot
\IfFileExists{footnotehyper.sty}{\usepackage{footnotehyper}}{\usepackage{footnote}}
\makesavenoteenv{longtable}
\setlength{\emergencystretch}{3em} % prevent overfull lines
\providecommand{\tightlist}{%
  \setlength{\itemsep}{0pt}\setlength{\parskip}{0pt}}
\emergencystretch=3em
\AtBeginDocument{\Urlmuskip=0mu plus 2mu}
\lstset{breaklines=true,basicstyle=\small\ttfamily,columns=fullflexible,keepspaces=true}
\setlength{\tabcolsep}{4pt}
\usepackage{iftex}
\ifXeTeX
  \usepackage{newunicodechar}
  \makeatletter
  \newcommand{\DeclareUnicodeCharacter}[2]{%
    \begingroup
    \uccode`\~="#1\relax
    \uppercase{\def\temp{~}}%
    \expandafter\endgroup
    \expandafter\newunicodechar\expandafter{\temp}{#2}%
  }
  \makeatother
\fi
\ifdefined\DeclareUnicodeCharacter
\DeclareUnicodeCharacter{2192}{\ensuremath{\to}}
\DeclareUnicodeCharacter{2190}{\ensuremath{\leftarrow}}
\DeclareUnicodeCharacter{2194}{\ensuremath{\leftrightarrow}}
\DeclareUnicodeCharacter{21D2}{\ensuremath{\Rightarrow}}
\DeclareUnicodeCharacter{21A6}{\ensuremath{\mapsto}}
\DeclareUnicodeCharacter{2212}{\ensuremath{-}}
\DeclareUnicodeCharacter{2264}{\ensuremath{\leq}}
\DeclareUnicodeCharacter{2265}{\ensuremath{\geq}}
\DeclareUnicodeCharacter{2260}{\ensuremath{\neq}}
\DeclareUnicodeCharacter{2113}{\ensuremath{\ell}}
\DeclareUnicodeCharacter{03A3}{\ensuremath{\Sigma}}
\DeclareUnicodeCharacter{03A6}{\ensuremath{\Phi}}
\DeclareUnicodeCharacter{0393}{\ensuremath{\Gamma}}
\DeclareUnicodeCharacter{221E}{\ensuremath{\infty}}
\DeclareUnicodeCharacter{2208}{\ensuremath{\in}}
\DeclareUnicodeCharacter{2209}{\ensuremath{\notin}}
\DeclareUnicodeCharacter{211A}{\ensuremath{\mathbb{Q}}}
\DeclareUnicodeCharacter{211D}{\ensuremath{\mathbb{R}}}
\DeclareUnicodeCharacter{2124}{\ensuremath{\mathbb{Z}}}
\DeclareUnicodeCharacter{2115}{\ensuremath{\mathbb{N}}}
\DeclareUnicodeCharacter{2200}{\ensuremath{\forall}}
\DeclareUnicodeCharacter{2203}{\ensuremath{\exists}}
\DeclareUnicodeCharacter{2286}{\ensuremath{\subseteq}}
\DeclareUnicodeCharacter{00D7}{\ensuremath{\times}}
\DeclareUnicodeCharacter{03B4}{\ensuremath{\delta}}
\DeclareUnicodeCharacter{2202}{\ensuremath{\partial}}
\DeclareUnicodeCharacter{03B2}{\ensuremath{\beta}}
\DeclareUnicodeCharacter{03B1}{\ensuremath{\alpha}}
\DeclareUnicodeCharacter{03C6}{\ensuremath{\phi}}
\DeclareUnicodeCharacter{03B3}{\ensuremath{\gamma}}
\DeclareUnicodeCharacter{00B9}{\ensuremath{^1}}
\DeclareUnicodeCharacter{00B2}{\ensuremath{^2}}
\DeclareUnicodeCharacter{2070}{\ensuremath{^0}}
\DeclareUnicodeCharacter{2080}{\ensuremath{_0}}
\DeclareUnicodeCharacter{2081}{\ensuremath{_1}}
\DeclareUnicodeCharacter{2082}{\ensuremath{_2}}
\fi
\usepackage{bookmark}
\IfFileExists{xurl.sty}{\usepackage{xurl}}{} % add URL line breaks if available
\urlstyle{same}
\hypersetup{
  pdftitle={PFP X: Exact Host Checking of Model Verdicts on Small Finite Algebraic Challenges --- Lattice Holonomy and Six Formal-Mathematics Pilots},
  pdfauthor={Jeremy H. Carroll},
  colorlinks=true,
  linkcolor={Maroon},
  filecolor={Maroon},
  citecolor={Blue},
  urlcolor={Blue},
  pdfcreator={LaTeX via pandoc}}

\title{PFP X: Exact Host Checking of Model Verdicts on Small Finite
Algebraic Challenges --- Lattice Holonomy and Six Formal-Mathematics
Pilots}
\author{Jeremy H. Carroll}
\date{2026-10-02 · Version 0.1.0 · Research methods paper ·
Author-review revision}

\begin{document}
\makeatletter
\begin{center}
{\LARGE \@title\par}
\vspace{0.8em}
{\normalsize \@author\par}
\vspace{0.5em}
{\small \@date\par}
\end{center}
\makeatother

\textbf{Code baseline:}
\nolinkurl{a623568eb0f853b968bc419bf06dfc6af256ab02}

\subsection{Abstract}\label{abstract}

A language model can state a verdict about a finite algebraic object. A
host program can check that verdict exactly. This paper supplies a
correct checker for one such task and reports what exact checking found
in two small pilots. For connections with values in a finite group on a
square lattice, the checker multiplies every plaquette in the
orientation fixed by the lattice, uses inverses on reversed edges, and
on a torus uses the products around the two non-contractible cycles to
separate flat connections that are gauge-trivial from those that are
not. Skipping plaquettes whose four links are the identity is an
optimization only; after a gauge transformation such plaquettes need not
exist. The checker replaces an incorrect function in the first release
of this package. On five retained cases from challenge C017 it confirms
the verdicts that the challenge harness expected, and it agrees with the
harness on all 15 model certificates that were retained in full. In five
runs (four prompting arms, one of them with two models) the harness
accepted 0, 0, 3, 3 and 4 of the 5 certificates; the arms differ in what
the prompt told the model, and n = 5. For six formal-mathematics
challenges (C045 to C050), three models each answered ten cases per
challenge; two challenges have two retained runs with different results.
We report the counts, what the paired checkers accepted, and the harness
limits, and we draw no comparison between models. These counts do not
measure unaided performance. Each prompt began with the case identifier,
and for 36 of the 60 cases in the six panels whose runner is retained
the identifier names the expected verdict or answer; in two of the
challenges an answer that refuses every case passes all ten. This
revision incorporates the earlier erratum. The first release's claims
about a five-stage prompt filter, a 50 of 50 benchmark and three effect
sizes are withdrawn, because no retained run supports them.

\subsection{1. Problem and design}\label{problem-and-design}

The Proposal Fidelity Protocol (PFP), developed in this collection under
the legacy name Symbol Game, separates a candidate answer from the host
check that decides whether the answer is accepted. This paper applies
that separation to model answers about small finite algebraic objects.
The model states a verdict and a certificate. A deterministic host
program recomputes the property from the problem data using only exact
operations: table lookups and integer arithmetic. The decision depends
on the certificate, not on the model's explanation.

Lattice holonomy is a useful test of this design. For a nonabelian
group, the product around a plaquette depends on the order of
multiplication and on the orientation, and an error in either gives a
confident wrong answer. Checking is cheap: one pass over the plaquettes
and, on a torus, two cycle products. The task therefore shows both
whether a host catches order errors and what a harness must tell a model
for its answer to be checkable at all.

The paper makes three contributions:

\begin{enumerate}
\def\labelenumi{\arabic{enumi}.}
\tightlist
\item
  a complete flatness check for finite-group connections on square
  lattices, with tests (Section 2);
\item
  an exact replay of the five retained C017 lattice cases and of every
  retained model certificate for them, with per-arm counts (Section 3);
\item
  a table of the retained formal-mathematics pilot runs for challenges
  C045 to C050, with what the paired checkers accepted, what the prompts
  disclosed and the limits the harness set (Section 4).
\end{enumerate}

No new model runs were made for this revision. Version 2.0.1 adds a
replay, with no model call, of answers that need no mathematics against
the C045 to C050 checkers as they were when the runs were recorded
(Section 4.2).

The first release of this paper (2026-09-28) also described a five-stage
prompt filter whose stages were named after immune mechanisms, and it
reported a 50 of 50 benchmark, a 40 per cent reduction in premature
rollbacks, an 83.5 per cent reduction in prompt tokens and a 50 per cent
accuracy gain. No run, runner or receipt for any of these is retained in
the repository, so they are withdrawn; the change log lists every
removed statement. In the first release's package the stages were simple
functions. ``MHC peptide decomposition'' split an edge list into
radius-one neighbourhoods. ``Thymic selection'' tested membership in a
supplied list. ``Kinetic proofreading'' compared two supplied fractions
with two thresholds. The ``Treg scratchpad'' check tested that four
labelled lines appear in order. The ``allosteric quorum'' counted how
many of five listed Unicode symbols occur in the prompt and compared the
count with three. No effect of any of them on model answers was
measured.

\subsection{2. A complete flatness
check}\label{a-complete-flatness-check}

\subsubsection{2.1 Definitions}\label{definitions}

A finite group \(G\) is given by named elements, a full multiplication
table and an inverse map. The package checks the identity law and the
inverse law for every element and associativity for every triple before
it uses a table, and it refuses a table that fails. It builds \(S_3\)
from the permutations of \(\{0,1,2\}\) and the quaternion group \(Q_8\)
from the unit quaternions \(\pm 1, \pm i, \pm j, \pm k\) with integer
coordinates.

An \(L \times M\) lattice has vertices \((r,c)\) with \(0 \le r < L\)
and \(0 \le c < M\). Its positive edges are \((r,c) \to (r,c+1)\) and
\((r,c) \to (r+1,c)\); on a torus the indices are taken modulo \(L\) and
\(M\). A torus needs \(L, M \ge 3\) so that each edge is named by its
two endpoints. A connection assigns an element \(U_{uv} \in G\) to each
positive edge, and a reversed edge carries the inverse,
\(U_{vu} = U_{uv}^{-1}\). The product along a path
\(v_0 \to v_1 \to \cdots \to v_k\) is
\[\operatorname{hol}(v_0 \to \cdots \to v_k) = U_{v_0 v_1} U_{v_1 v_2} \cdots U_{v_{k-1} v_k}.\]
The plaquette \(P_{r,c}\) is listed as \(v_0 = (r,c)\),
\(v_1 = (r+1,c)\), \(v_2 = (r+1,c+1)\), \(v_3 = (r,c+1)\), and its
product is
\[\operatorname{hol}(P_{r,c}) = U_{v_0 v_1} U_{v_1 v_2} U_{v_2 v_3} U_{v_3 v_0}.\]
The last two factors traverse positive edges backwards, so they are the
inverses of the stored links. A connection is flat when every plaquette
product is the identity \(e\).

A gauge \(g\) assigns an element to each vertex and acts by
\(U'_{uv} = g_u U_{uv} g_v^{-1}\). The product around a closed path
based at \(v_0\) becomes \(g_{v_0} \operatorname{hol} g_{v_0}^{-1}\).
Flatness is therefore unchanged by a gauge transformation. On a torus,
let \(A\) be the product around
\((0,0) \to (0,1) \to \cdots \to (0,M-1) \to (0,0)\) and \(B\) the
product around \((0,0) \to (1,0) \to \cdots \to (L-1,0) \to (0,0)\). A
gauge transformation conjugates \(A\) and \(B\) by the same element.

We use four standard facts {[}1, 2, 3{]}. Reversing a closed path
inverts its product. Starting the same closed path at \(v_1\) instead of
\(v_0\) conjugates its product by \(U_{v_0 v_1}^{-1}\). If a connection
is flat, the product along a closed path depends only on the path's
homotopy class, so on an open lattice, which is simply connected, a
gauge built along a spanning tree makes every link the identity. On the
torus, whose fundamental group \(\mathbb{Z}^2\) is abelian, a flat
connection has commuting \(A\) and \(B\), and it is gauge-trivial
exactly when \(A = B = e\).

\subsubsection{2.2 The check}\label{the-check}

The checker in \passthrough{\lstinline!code/holonomy.py!} returns one of
three verdicts.

\begin{enumerate}
\def\labelenumi{\arabic{enumi}.}
\tightlist
\item
  It reads the group table and refuses it unless it is a group. It reads
  the links and refuses them unless every edge is listed in both
  directions with \(U_{vu} = U_{uv}^{-1}\) and no other edge is listed.
\item
  It computes every plaquette product in the order of Section 2.1. If
  any is not \(e\), it returns \passthrough{\lstinline!NOT\_FLAT!} with
  the first such plaquette in row-major order and its product.
\item
  On an open lattice it returns \passthrough{\lstinline!FLAT\_TRIVIAL!}
  with the gauge \(g_{(0,0)} = e\), \(g_v = g_u U_{uv}\) along a
  breadth-first spanning tree, after verifying that this gauge makes
  every link the identity.
\item
  On a torus it computes \(A\) and \(B\) and stops with an error if they
  do not commute, which cannot happen for a flat connection. If
  \(A = B = e\) it returns \passthrough{\lstinline!FLAT\_TRIVIAL!} with
  the verified tree gauge; otherwise it returns
  \passthrough{\lstinline!FLAT\_NONTRIVIAL!} with \(A\) and \(B\).
\end{enumerate}

\passthrough{\lstinline!FLAT\_TRIVIAL!} is the verdict that challenge
C017 calls \passthrough{\lstinline!PURE\_GAUGE!}.

A separate function checks a stated verdict and its certificate. For
\passthrough{\lstinline!NOT\_FLAT!} the four listed vertices must bound
a plaquette, in either orientation and from any starting vertex, and the
stated product must equal the product in the listed order and differ
from \(e\). For \passthrough{\lstinline!FLAT\_NONTRIVIAL!} every
plaquette must be flat and at least one stated closed loop must have the
stated, non-identity product. In a flat connection every contractible
loop has product \(e\), so such a loop is non-contractible. For
\passthrough{\lstinline!FLAT\_TRIVIAL!} the stated gauge must make every
link the identity. Any other answer, including an answer that could not
be parsed, is rejected with a reason.

\subsubsection{2.3 Remark: skipping plaquettes with four identity
links}\label{remark-skipping-plaquettes-with-four-identity-links}

If the four links of a plaquette are the identity, its product is the
identity. This uses only the identity law and \(e^{-1} = e\); the
package records the one-line proof in Lean
(\nolinkurl{identity_links_give_identity_product} in
\mbox{\nolinkurl{lean4/LatticeDecimation.lean}}). The first release
called this observation a theorem. It is a direct consequence of the
group axioms.

The checker may skip such plaquettes. It then performs fewer
multiplications and reaches the same result, which the tests confirm on
every connection they build. The saving depends on the gauge. If
\(U = e\) and \(g\) is a gauge, then \(U'_{uv} = g_u g_v^{-1}\) is the
identity only when \(g_u = g_v\). A gauge that differs on neighbouring
vertices therefore leaves no plaquette with four identity links, while
the connection stays flat and gauge-trivial. One retained C017 case,
\passthrough{\lstinline!base\_02!}, has this form: 28 of its 31 links
are not the identity, none of its 12 plaquettes can be skipped, and it
is gauge-trivial. Skipping is an optimization. It cannot replace the
products of the remaining plaquettes, and it cannot show flatness on its
own.

\subsubsection{2.4 Tests}\label{tests}

The tests in \nolinkurl{code/tests/test_holonomy_checker.py} check the
following on \(S_3\) and \(Q_8\):

\begin{itemize}
\tightlist
\item
  the trivial connection is \passthrough{\lstinline!FLAT\_TRIVIAL!} on
  open and periodic lattices;
\item
  a connection that is a gauge transformation of the trivial one, with
  no plaquette of four identity links, is
  \passthrough{\lstinline!FLAT\_TRIVIAL!};
\item
  a torus connection with a constant non-identity product around one
  cycle is \passthrough{\lstinline!FLAT\_NONTRIVIAL!};
\item
  a single non-identity link away from the boundary gives
  \passthrough{\lstinline!NOT\_FLAT!} with exactly two non-flat
  plaquettes;
\item
  for every pair \((a,b) \in G \times G\), the \(3 \times 4\) torus
  connection whose links from the last column to the first carry \(a\),
  whose links from the last row to the first carry \(b\), and whose
  other links are \(e\), is flat exactly when \(ab = ba\); when it is
  flat its cycle products are \((a, b)\), or their conjugates by one
  element after a gauge transformation, and when it is not flat its only
  non-flat plaquette has product \(b a b^{-1} a^{-1}\);
\item
  on seeded random connections, reversing a plaquette inverts its
  product, changing the starting vertex conjugates it, and some products
  differ from their inverses;
\item
  skipping plaquettes with four identity links never changes the result;
\item
  malformed group tables, lattices and link lists are refused, and the
  certificate check accepts a correct plaquette in either orientation
  and rejects a wrong product, a flat plaquette, a four-cycle that is
  not a plaquette and a contractible loop.
\end{itemize}

The seam family is small, but it is exhaustive over \(G \times G\), and
in it a flat connection with non-commuting cycle products never
occurred, as the theory requires.

\subsection{3. The C017 pilot}\label{the-c017-pilot}

\subsubsection{3.1 What was run}\label{what-was-run}

Challenge C017 asks for \passthrough{\lstinline!NOT\_FLAT!},
\passthrough{\lstinline!PURE\_GAUGE!} or
\passthrough{\lstinline!FLAT\_NONTRIVIAL!} verdicts with certificates
for connections with values in \(S_3\), \(Q_8\) and
\(\mathrm{SL}(2, \mathbb{F}_3)\) on square lattices. Its specification
calls for open-weight models of at most 70B parameters and 36 hidden
base cases. The retained runs are smaller pilots. They used five cases
from the challenge's judge packet (Table 1), one reply per case, and one
model outside that specification. The five case files are copied into
\passthrough{\lstinline!code/data/c017/!} with their SHA-256 digests and
source paths.

Each arm sent one prompt per case and parsed a JSON object from the
reply. Two checkers in the harness scored the certificate, and a case
passed only when both accepted it.

\begin{itemize}
\tightlist
\item
  \textbf{Arm 1.} The lattice, the group table and inverses, all
  directed links, the four verdict definitions and the JSON format.
  Model \nolinkurl{meta-llama/llama-3.3-70b-instruct}; temperature 0.1;
  at most 2,000 completion tokens.
\item
  \textbf{Arm 2.} Arm 1 plus four written constraints: multiply in
  traversal order; test the product, not a sign or determinant; on a
  torus a flat connection need not be gauge-trivial; refuse inverse
  mismatches. Same model and settings.
\item
  \textbf{Arm 5.} The Arm 2 constraints plus the full list of plaquettes
  in the vertex order that the checker accepts, an instruction to
  multiply step by step by table lookup, the two cycle paths on a torus
  and a decision procedure. Run with
  \nolinkurl{meta-llama/llama-3.3-70b-instruct} (temperature 0.05, at
  most 3,500 completion tokens) and with
  \nolinkurl{stealth/space-bunny-alpha} (temperature 0.05; the runner's
  default limit is 6,000 completion tokens, can be changed on the
  command line, and is not recorded in the report).
\item
  \textbf{Arm 6.} Arm 5, except that the harness listed only the
  plaquettes that contain a non-identity link and stated that the others
  were flat. Run with \nolinkurl{stealth/space-bunny-alpha} (temperature
  0.05; runner default limit 4,000 completion tokens).
\end{itemize}

No C017 runs are retained for Arms 3 and 4. The model
\nolinkurl{stealth/space-bunny-alpha} is a new free ``stealth'' model
offered on OpenRouter as of 2026-09-28; the provider does not disclose
its identity, size or training. Results with it are exploratory. They
fall outside the challenge protocol, SG-CP v1, whose §2 admits only
open-weight models of at most 70B parameters
(\nolinkurl{projects/mssn_sn_symgam/challenges/01_challenge_protocol.md}),
and they may not be reproducible once the stealth period ends. The
reports record costs computed by the harness from configured per-token
prices, not provider bills, so we do not report costs.

Table 1. The five retained C017 cases and the host verdicts computed by
the package. ``Skippable'' counts plaquettes whose four links are the
identity. The case files name the elements of \(S_3\) as
\passthrough{\lstinline!e!}, \passthrough{\lstinline!r!},
\passthrough{\lstinline!r2!}, \passthrough{\lstinline!s!},
\passthrough{\lstinline!sr!}, \passthrough{\lstinline!sr2!}, where
\passthrough{\lstinline!r2!} is \(r^2 = r^{-1}\), and the elements of
\(Q_8\) as \passthrough{\lstinline!e!}, \passthrough{\lstinline!-e!},
\passthrough{\lstinline!i!}, \passthrough{\lstinline!-i!},
\passthrough{\lstinline!j!}, \passthrough{\lstinline!-j!},
\passthrough{\lstinline!k!}, \passthrough{\lstinline!-k!}. The case
\passthrough{\lstinline!pois\_01!} comes from the challenge's poisoned
set: its non-flat product is a 3-cycle, which has sign \(+1\), so a test
of the sign alone would call the connection flat.

\begin{longtable}[]{@{}
  >{\raggedright\arraybackslash}p{(\linewidth - 8\tabcolsep) * \real{0.2000}}
  >{\raggedright\arraybackslash}p{(\linewidth - 8\tabcolsep) * \real{0.2000}}
  >{\raggedright\arraybackslash}p{(\linewidth - 8\tabcolsep) * \real{0.2000}}
  >{\raggedright\arraybackslash}p{(\linewidth - 8\tabcolsep) * \real{0.2000}}
  >{\raggedright\arraybackslash}p{(\linewidth - 8\tabcolsep) * \real{0.2000}}@{}}
\toprule\noalign{}
\begin{minipage}[b]{\linewidth}\raggedright
Case
\end{minipage} & \begin{minipage}[b]{\linewidth}\raggedright
Group and lattice
\end{minipage} & \begin{minipage}[b]{\linewidth}\raggedright
Non-identity links
\end{minipage} & \begin{minipage}[b]{\linewidth}\raggedright
Plaquettes (skippable)
\end{minipage} & \begin{minipage}[b]{\linewidth}\raggedright
Host verdict
\end{minipage} \\
\midrule\noalign{}
\endhead
\bottomrule\noalign{}
\endlastfoot
\passthrough{\lstinline!pois\_01!} & \(S_3\), \(3 \times 3\) open & 1 of
12 & 4 (3) & \passthrough{\lstinline!NOT\_FLAT!}, \(P_{0,0}\) has
product \passthrough{\lstinline!r2!} \\
\passthrough{\lstinline!base\_01!} & \(S_3\), \(3 \times 4\) open & 1 of
17 & 6 (5) & \passthrough{\lstinline!NOT\_FLAT!}, \(P_{0,1}\) has
product \passthrough{\lstinline!s!} \\
\passthrough{\lstinline!base\_02!} & \(S_3\), \(4 \times 5\) open & 28
of 31 & 12 (0) & \passthrough{\lstinline!FLAT\_TRIVIAL!} \\
\passthrough{\lstinline!base\_03!} & \(S_3\), \(5 \times 6\) torus & 11
of 60 & 30 (20) & \passthrough{\lstinline!FLAT\_NONTRIVIAL!}, \(A\) =
\passthrough{\lstinline!r!}, \(B\) = \passthrough{\lstinline!r2!} \\
\passthrough{\lstinline!base\_13!} & \(Q_8\), \(4 \times 5\) open & 1 of
31 & 12 (11) & \passthrough{\lstinline!NOT\_FLAT!}, \(P_{0,0}\) has
product \passthrough{\lstinline!-i!} \\
\end{longtable}

The host verdicts agree with the verdicts that the harness expected for
all five cases. The \(S_3\) table in the case files is isomorphic to the
permutation group (the tests find all six isomorphisms), and the \(Q_8\)
table equals the package's quaternion table.

\subsubsection{3.2 Outcomes}\label{outcomes}

Table 2 gives, for each arm and case, what the model stated and whether
the harness accepted the certificate. ``70B'' denotes
\nolinkurl{meta-llama/llama-3.3-70b-instruct} and ``stealth'' denotes
\nolinkurl{stealth/space-bunny-alpha}. ``Reversed order'' means that the
answer named the correct non-flat plaquette with its vertices in the
opposite orientation.

Table 2. Harness outcomes for the five C017 cases (one reply per cell).

\begin{longtable}[]{@{}
  >{\raggedright\arraybackslash}p{(\linewidth - 10\tabcolsep) * \real{0.1667}}
  >{\raggedright\arraybackslash}p{(\linewidth - 10\tabcolsep) * \real{0.1667}}
  >{\raggedright\arraybackslash}p{(\linewidth - 10\tabcolsep) * \real{0.1667}}
  >{\raggedright\arraybackslash}p{(\linewidth - 10\tabcolsep) * \real{0.1667}}
  >{\raggedright\arraybackslash}p{(\linewidth - 10\tabcolsep) * \real{0.1667}}
  >{\raggedright\arraybackslash}p{(\linewidth - 10\tabcolsep) * \real{0.1667}}@{}}
\toprule\noalign{}
\begin{minipage}[b]{\linewidth}\raggedright
Case
\end{minipage} & \begin{minipage}[b]{\linewidth}\raggedright
Arm 1, 70B
\end{minipage} & \begin{minipage}[b]{\linewidth}\raggedright
Arm 2, 70B
\end{minipage} & \begin{minipage}[b]{\linewidth}\raggedright
Arm 5, 70B
\end{minipage} & \begin{minipage}[b]{\linewidth}\raggedright
Arm 5, stealth
\end{minipage} & \begin{minipage}[b]{\linewidth}\raggedright
Arm 6, stealth
\end{minipage} \\
\midrule\noalign{}
\endhead
\bottomrule\noalign{}
\endlastfoot
\passthrough{\lstinline!pois\_01!} & reversed order & reversed order &
accepted & accepted & accepted \\
\passthrough{\lstinline!base\_01!} & reversed order & stated
\passthrough{\lstinline!PURE\_GAUGE!} & accepted & accepted &
accepted \\
\passthrough{\lstinline!base\_02!} & no parsable JSON & stated
\passthrough{\lstinline!NOT\_FLAT!} & flat plaquette stated as
\passthrough{\lstinline!r2!} & flat plaquette stated as
\passthrough{\lstinline!r!} & flat plaquette stated as
\passthrough{\lstinline!r!} \\
\passthrough{\lstinline!base\_03!} & loop product misstated & loop
product misstated & no parsable JSON & reply cut off & accepted \\
\passthrough{\lstinline!base\_13!} & reversed order & no parsable JSON &
accepted & accepted & accepted \\
Accepted & 0 of 5 & 0 of 5 & 3 of 5 & 3 of 5 & 4 of 5 \\
Verdict label equal to host & 4 of 5 & 2 of 5 & 3 of 5 & 3 of 5 & 4 of
5 \\
\end{longtable}

For the 15 certificates retained in full (Arms 5 and 6), the package's
certificate check agrees with the harness in every case: 10 accepted and
5 rejected. For Arms 1 and 2 the reports retain only the stated verdict
and the checkers' reasons, not the certificates.

\textbf{Vertex order.} The Arm 1 and Arm 2 prompts defined a plaquette
product as \(U_{01} U_{12} U_{23} U_{30}\) but did not say in which
order the checker lists a plaquette's vertices. The checker accepted
only \((r,c), (r+1,c), (r+1,c+1),
(r,c+1)\). Three Arm 1 answers named the correct non-flat plaquette in
the order \((r,c), (r,c+1), (r+1,c+1), (r+1,c)\) and were rejected on
that ground before any product was compared. The reports do not retain
the stated products, so we cannot say whether they were right for the
reversed order. Arm 1 therefore scored 0 of 5, although its verdict
label matched the host in 4 of 5 cases. Two Arm 2 answers were rejected
on the same ground; one of them also had the wrong verdict. The Arm 5
prompt listed the accepted order, and the three open non-flat cases were
accepted for both models. Arm 5 also added the step-by-step instruction
and the cycle paths, so these runs do not separate the effect of stating
the convention from the other additions. The certificate check of
Section 2 accepts either orientation.

\textbf{The gauge-trivial case.} \passthrough{\lstinline!base\_02!} is
flat and gauge-trivial, and 28 of its 31 links are not the identity. All
three Arm 5 and Arm 6 answers for it stated
\passthrough{\lstinline!NOT\_FLAT!} with a plaquette whose product is
\(e\), stating the product as \passthrough{\lstinline!r2!} or
\passthrough{\lstinline!r!}; each was rejected because the stated
product disagrees with the table. The first release also stated that an
instruction to build a gauge along a spanning tree biased models towards
\passthrough{\lstinline!PURE\_GAUGE!}. No retained run supports this,
and the statement is withdrawn. The only retained
\passthrough{\lstinline!PURE\_GAUGE!} answer is the Arm 2 answer for
\passthrough{\lstinline!base\_01!}, a non-flat case, whose stated gauge
fails on the edge \((0,1) \to (0,2)\).

\textbf{The torus case.} In \passthrough{\lstinline!base\_03!}, 20 of
the 30 plaquettes have four identity links, and each of the other 10
contains at least one of the 11 non-identity links. With Arm 5, the
reply of \nolinkurl{stealth/space-bunny-alpha} stopped at 8,000
completion tokens with its JSON unfinished; the retained start of the
reply already carried the label
\passthrough{\lstinline!FLAT\_NONTRIVIAL!}. With Arm 6, where the
harness listed only the 10 plaquettes with a non-identity link, the same
model returned all 10 products (each \(e\)), \(A\) =
\passthrough{\lstinline!r!}, \(B\) = \passthrough{\lstinline!r2!} and
the label \passthrough{\lstinline!FLAT\_NONTRIVIAL!} in 2,810 completion
tokens, and both checkers accepted the answer. This is one case with one
reply per arm. It shows that the shorter prompt fitted within the limit
on this case; it does not measure a general saving. The Arm 6 listing
relies on identity links, which a gauge transformation can remove
(Section 2.3). The first release stated that 24 plaquettes were skipped
and that the model evaluated only the two cycles. Recomputing the Arm 6
plaquette list from the case file gives 20 skipped plaquettes, and the
retained reply evaluates the other 10 plaquettes and the two cycles.

\textbf{The harness checkers.} The two checkers in the C017 harness are
not independent implementations. Both import the plaquette list and the
path product from the challenge's reference solver. The check in this
package shares no code with them, and on the retained certificates it
reaches the same decisions.

\subsubsection{3.3 Data sources for Section
3}\label{data-sources-for-section-3}

All paths are relative to the repository root.

\begin{itemize}
\tightlist
\item
  Challenge specification:
  \nolinkurl{projects/mssn_sn_symgam/challenges/b_cohomology/c017_nonabelian_lattice_holonomy.md}.
\item
  Arms 1 and 2:
  \nolinkurl{.../c017_nonabelian_lattice_holonomy/results/immune_transfer_report.json},
  added in commit \nolinkurl{7e639a0be573f6affc4d84a2aacb5573da92bc79};
  runner \nolinkurl{eval_epistemic_immune_transfer.py} at that commit.
\item
  Arm 5 with \nolinkurl{meta-llama/llama-3.3-70b-instruct}:
  \nolinkurl{results/arm5_hybrid_report.json} as committed in
  \nolinkurl{9e45ee501e6a48086aaabd2381632eef6f7ed7a2}. The file at that
  path was later overwritten by a one-case run, so this report is
  retained only in the repository history.
\item
  Arm 5 with \nolinkurl{stealth/space-bunny-alpha}:
  \nolinkurl{results/arm5_stealth_space_bunny_alpha_report.json}, commit
  \nolinkurl{5327cfcb7ee3622e79f7d0cd884dfa37ef4eef56}.
\item
  Arm 6: \nolinkurl{results/arm6_stealth_space_bunny_alpha_report.json},
  commit \nolinkurl{eb5fc971ae44320ff8183f09343d547f830d9d56}.
\end{itemize}

Each report's SHA-256 digest equals the digest recorded in its receipt
file in the same folder. The package file
\nolinkurl{code/data/c017/model_claims.json} holds the verdicts,
certificates, checker reasons and completion lengths extracted from
these reports.

\subsection{4. Six formal-mathematics
pilots}\label{six-formal-mathematics-pilots}

\subsubsection{4.1 What was run}\label{what-was-run-1}

For each of challenges C045 to C050, a runner script defines a ten-case
panel, separate from the challenge's 64 hidden cases (for C049 panel B
the script is not retained; see below), and sends each case once to
three models through OpenRouter at temperature 0:
\nolinkurl{stealth/space-bunny-alpha} (the undisclosed stealth model
described in Section 3.1), \nolinkurl{meta-llama/llama-3.3-70b-instruct}
and \nolinkurl{meta-llama/llama-3.1-8b-instruct}. The reports do not
record which provider served each model. A case passed only when both of
the challenge's checkers, an exact checker and a second checker,
accepted the answer.

The harness set the completion limits. In C045 and C046 the limit was
4,096 tokens for \nolinkurl{stealth/space-bunny-alpha} and 2,048 for the
other two models. In C047 and C049 panel A it was 4,096 tokens and in
C048 and C050 2,048 tokens for all three. The limit for C049 panel B is
not recorded.

Table 3. Cases passed, out of 10 per cell, one reply per case.

\begin{longtable}[]{@{}
  >{\raggedright\arraybackslash}p{(\linewidth - 8\tabcolsep) * \real{0.2000}}
  >{\raggedright\arraybackslash}p{(\linewidth - 8\tabcolsep) * \real{0.2000}}
  >{\raggedright\arraybackslash}p{(\linewidth - 8\tabcolsep) * \real{0.2000}}
  >{\raggedright\arraybackslash}p{(\linewidth - 8\tabcolsep) * \real{0.2000}}
  >{\raggedright\arraybackslash}p{(\linewidth - 8\tabcolsep) * \real{0.2000}}@{}}
\toprule\noalign{}
\begin{minipage}[b]{\linewidth}\raggedright
Challenge
\end{minipage} & \begin{minipage}[b]{\linewidth}\raggedright
Task
\end{minipage} & \begin{minipage}[b]{\linewidth}\raggedright
\nolinkurl{stealth/space-bunny-alpha}
\end{minipage} & \begin{minipage}[b]{\linewidth}\raggedright
llama-3.3-70b-instruct
\end{minipage} & \begin{minipage}[b]{\linewidth}\raggedright
llama-3.1-8b-instruct
\end{minipage} \\
\midrule\noalign{}
\endhead
\bottomrule\noalign{}
\endlastfoot
C045 & knot invariants from planar diagrams & 2 & 8 & 3 \\
C046 & Knuth--Bendix completion and normal forms {[}4{]} & 4 & 7 & 8 \\
C047 & linear recurrences or certified absence & 4 & 2 & 0 \\
C048 & rational point or local obstruction for a ternary quadratic form
& 10 & 8 & 5 \\
C049, panel A & limits, colimits and adjunctions in finite categories &
5 & 6 & 4 \\
C049, panel B & the same topics, a different set of ten cases & 9 & 6 &
6 \\
C050, run 1 & equational rewriting proof or two-element countermodel & 7
& 4 & 1 \\
C050, run 2 & the same ten cases, run again & 6 & 3 & 1 \\
\end{longtable}

C049 has two retained runs on different ten-case panels with the same
three models, committed together on 2026-09-28. Panel A was produced by
the runner script retained in the repository; the script that produced
panel B is not retained. For \nolinkurl{stealth/space-bunny-alpha} the
two panels give 5 and 9 passes. C050 was run twice with the same
committed runner and the same ten cases; the recorded start times are 98
seconds apart. The second report overwrote the first, which is retained
in the repository history. For \nolinkurl{stealth/space-bunny-alpha},
five of the ten cases changed outcome between the two runs, two of them
because one run received an API error instead of a reply. The first
release reported only C049 panel B and C050 run 2.

\textbf{What the prompts disclosed.} Each runner sent a case as the line
\passthrough{\lstinline!Case ID: <id>!} followed by the case data (lines
219, 270, 221, 163, 264 and 236 of the C045 to C050 runners, unchanged
from the commits that recorded the runs to the code baseline). The
challenge authors wrote the identifiers, and many name the expected
result. \nolinkurl{c045_m02_trefoil_identify} asks for the
identification of a knot that its identifier names.
\nolinkurl{c045_m09_refusal_corrupt_edge} and
\nolinkurl{c046_m08_refusal_commutativity} name a refusal and its
reason. \nolinkurl{c046_m05_check_non_confluent_group} names the
confluence verdict. \passthrough{\lstinline!c047\_m07\_prand\_decoy!}
and \passthrough{\lstinline!c047\_m08\_order9\_decoy!} mark the cases
with no recurrence within the bounds.
\passthrough{\lstinline!c048\_m01\_solv\_11\_m2!} and
\passthrough{\lstinline!c048\_m02\_solv\_pyth!} mark solvable forms, and
\passthrough{\lstinline!c048\_m05\_obstr\_inf\_111!} names both the
obstruction and its place. \nolinkurl{c049_m05_no_product_disc2} names
the verdict that no product exists. The C045, C046, C049 and C050
runners also sent the whole case record, including the
\passthrough{\lstinline!note!} of the poisoned cases C049 m10 and C050
m09, whose identifiers call the note bait or fake; the C047 and C048
runners sent only the terms or coefficients and the bounds. Table 4
counts the identifiers of the six panels whose runner is retained. Which
identifiers name the result is one reader's classification, the
author's; it is published identifier by identifier with the replay
listed in Section 4.4, so that another reader can check or dispute each
case.

Table 4. What the case identifiers name, in the author's classification,
and the most passes that an answer needing no mathematics obtains under
the same rule (both checkers accept), out of 10 per panel. The answers
tried were an empty object, a bare abstention, a refusal or out-of-scope
answer with each reason the prompt lists, and a constant verdict where
the output format allows one (Section 4.2).

\begin{longtable}[]{@{}
  >{\raggedright\arraybackslash}p{(\linewidth - 8\tabcolsep) * \real{0.1579}}
  >{\raggedleft\arraybackslash}p{(\linewidth - 8\tabcolsep) * \real{0.2105}}
  >{\raggedleft\arraybackslash}p{(\linewidth - 8\tabcolsep) * \real{0.2105}}
  >{\raggedleft\arraybackslash}p{(\linewidth - 8\tabcolsep) * \real{0.2105}}
  >{\raggedleft\arraybackslash}p{(\linewidth - 8\tabcolsep) * \real{0.2105}}@{}}
\toprule\noalign{}
\begin{minipage}[b]{\linewidth}\raggedright
Challenge
\end{minipage} & \begin{minipage}[b]{\linewidth}\raggedleft
Names the verdict or answer
\end{minipage} & \begin{minipage}[b]{\linewidth}\raggedleft
Names the object
\end{minipage} & \begin{minipage}[b]{\linewidth}\raggedleft
Neither
\end{minipage} & \begin{minipage}[b]{\linewidth}\raggedleft
Best answer without mathematics
\end{minipage} \\
\midrule\noalign{}
\endhead
\bottomrule\noalign{}
\endlastfoot
C045 & 7 & 3 & 0 & 10 \\
C046 & 5 & 3 & 2 & 10 \\
C047 & 4 & 6 & 0 & 0 \\
C048 & 10 & 0 & 0 & 3 \\
C049, panel A & 5 & 5 & 0 & 2 \\
C050 & 5 & 5 & 0 & 1 \\
\end{longtable}

In C045 and C046 the best such answer is a refusal of every case: both
checkers accept a refusal with any of the listed reasons on every case,
so it passes 10 of 10, more than any model in Table 3. In C048 a
constant local obstruction at the place 2 passes 3 cases, and in C049 a
constant ``does not exist'' passes 2. The identifiers of the C049 panel
B report have the same form (for example
\nolinkurl{c049_m07_no_product_discrete}), but its runner is not
retained, so what its prompts contained cannot be checked and its cases
cannot be replayed.

The counts in Table 3 are therefore counts under leaked labels. These
counts describe responses under the disclosed prompts and checker rules;
they do not estimate or bound unaided performance. In C045 and C046 they
do not separate any model from an answer that refuses every case. The
models' own passes there include refusals only on the cases whose
identifiers name a refusal (C045 m09 and m10; C046 m08 to m10), and the
checkers accepted those refusals without comparing the stated reason
with the input.

We do not add counts across challenges and we make no comparison between
models. The tasks differ, the panels were written by the challenge
authors, each cell has n = 10 and one reply per case, and no statistical
test was run. The first release's totals over 80 cases and its
conclusion that absence proofs remain a capability of symbolic
algorithms alone are withdrawn; the evidence is 10 cases per cell on one
day.

\subsubsection{4.2 Observations}\label{observations}

\textbf{Harness limits.} In C045 all 8 failures of
\nolinkurl{stealth/space-bunny-alpha} returned no answer object, with
recorded token counts (prompt and completion together) of 4,751 to
4,845, consistent with the reply reaching the 4,096-token completion
limit. In C046, 5 of its 6 failures were of the same kind, at 4,624 to
4,747 tokens. The reports do not record prompt and completion counts
separately. These cases most likely show replies that did not fit the
harness limit; they do not measure what the model could compute.

\textbf{The second checkers.} For those empty answers in C045 and C046
the second checker reported acceptance and only the exact checker
rejected them. In C047 case m06, whose terms after the first are the
Fibonacci numbers, \nolinkurl{stealth/space-bunny-alpha} returned a
certificate that no recurrence exists within the bounds; the second
checker accepted it, and the exact checker rejected it because its minor
used rows below the smallest index that the exact checker allows.

A replay of answers that need no mathematics, against the checkers as
they were when the runs were recorded, shows how far this goes. It
imports each challenge's two checkers from those commits, reads the
panel from the runner without running it, and calls no model. The second
checkers of C045 and C046 accept the empty object
\passthrough{\lstinline!\{\}!}, a bare abstention, and any refusal that
carries a reason and a witness, on all ten cases. The second checker of
C047 accepts a recurrence with an empty coefficient list, and an absence
certificate with an empty list of minor rows, on 9 of the 10 cases,
because the determinant of an empty matrix is 1; on the tenth, whose
terms include a non-integer, it raises an error. The second checker of
C048 accepts any out-of-scope answer on all ten cases, and that of C049
accepts a constant ``does not exist'' on 8 of 10. A pass required both
checkers, so these defects do not change Table 3. They do mean that in
these challenges the pair of checkers is no stronger than the exact
checker alone, and in C045 and C046 the exact checker accepts any
refusal that carries a reason (Table 4).

The challenge harness was corrected after these runs, in commits
\passthrough{\lstinline!e6fbc15cf!} and
\passthrough{\lstinline!e2f7a3500!}. The C045 and C046 second checkers
now recompute the answer independently, so a refusal of every case,
\passthrough{\lstinline!\{\}!} and an answer without content fail them;
the C047 second checker no longer accepts an empty coefficient list or
an empty list of minor rows; the C048 second checker accepts an
out-of-scope answer only when a coefficient is not an integer; and the
C043 to C050 runners replace the descriptive identifiers with opaque
tokens drawn for each run. The C045 and C046 exact checkers still accept
a refusal with any reason, so the pair now depends on the second checker
for that case. None of this changes Table 3, which was produced with the
checkers of the run commits.

\textbf{Absence certificates.} C047 certifies that no recurrence exists
within the bounds by a nonsingular minor, whose determinant the exact
checker computes by fraction-free elimination {[}5{]}. On the three
panel cases that the runner marks as having no recurrence within the
bounds (m07, m08, m10), no model returned a certificate that both
checkers accepted. \nolinkurl{stealth/space-bunny-alpha} returned no
answer content for all three; the other two models stated recurrences or
malformed certificates.

\textbf{Other failures.} In each C050 run, one failure of
\nolinkurl{stealth/space-bunny-alpha} was an API error with no reply. In
C048, the two failures of
\passthrough{\lstinline!llama-3.3-70b-instruct!} were integer points
that do not satisfy the form's equation.

\subsubsection{4.3 What the checkers and mutation suites
show}\label{what-the-checkers-and-mutation-suites-show}

Each challenge's judge packet retains 64 hidden cases and 48 mutated
certificates for testing the checkers, 512 and 384 in total over C043 to
C050. Scoring receipts are retained for C043 and C044 only. Each records
the reference solver passing all 64 hidden cases, the two checkers
agreeing on all of them, and all 48 mutations rejected. No scoring
receipt is retained for C045 to C050, so the first release's statement
that 512 cases were verified and 384 mutations rejected is narrowed to
these 128 cases and 96 mutations. The receipts concern the reference
solvers and checkers, not the models.

Challenges C043 and C044 were also run with the same three models. The
C043 runner wrote its report to a directory that the repository's ignore
rules exclude. The C044 runner wrote its report outside the repository,
to a directory it reaches six levels above its own file. Neither report
is under version control, and the counts appear only in narrative
findings documents, so we do not report them. They would not measure
identification in any case: the C044 runner sent each example file
whole, including its \passthrough{\lstinline!note!}, which names the
group or the expected verdict (for example ``Quaternion group Q8'', ``C4
against Klein four group V4'' and ``S5 has order 120 which exceeds max
allowed order 64''). The C043 prompts were built from the case fields
and carried no note. A file with the C044 report's name exists outside
the repository (SHA-256
\nolinkurl{5fecd63f08dc1b5b2094fe6b722c236f284567b52a9a140f62d4825e63e8648b});
it was not under version control when the run was recorded, and we do
not report its counts. In the corrected harness the C044 runner no
longer sends the \passthrough{\lstinline!note!} or reads a credentials
file, and it writes its report into the repository's run store.

\subsubsection{4.4 Data sources for Section
4}\label{data-sources-for-section-4}

All paths are relative to
\nolinkurl{projects/mssn_sn_symgam/challenges/e_formal_math/}.

\begin{itemize}
\tightlist
\item
  C045 run report:
  \nolinkurl{c045_knot_polynomials_from_pd/empirical_results/c045_space_bunny_matrix_report.json},
  commit \nolinkurl{ea2b55d841bdbe283f2f6ac77abe8a1677cbf0f8}.
\item
  C046, C047 and C048 run reports: the files with the same pattern in
  their challenge folders, commit
  \nolinkurl{46bd1dec156260f073b77caa07d8a453f8efad30}.
\item
  C049 panel B:
  \nolinkurl{c049_finite_category_universal_constructions/empirical_results/c049_space_bunny_matrix_report.json};
  panel A: the three files \passthrough{\lstinline!results\_*.json!} in
  the same folder; both in commit
  \nolinkurl{46bd1dec156260f073b77caa07d8a453f8efad30}.
\item
  C050 run 1:
  \nolinkurl{c050_equational_theorem_proving/empirical_results/c050_space_bunny_matrix_report.json}
  as committed in \nolinkurl{46bd1dec156260f073b77caa07d8a453f8efad30};
  run 2: the same path as committed in
  \nolinkurl{5a256c6e8d088c612332362a8abc53138806ac7a}.
\item
  Runner scripts, which set the panels, the prompts and the completion
  limits: \passthrough{\lstinline!run\_c045\_matrix.py!} to
  \passthrough{\lstinline!run\_c050\_matrix.py!} in the challenge
  folders.
\item
  The replay behind Table 4 and Section 4.2 (paths relative to the
  repository root): the collection's script
  \nolinkurl{atlas/coda_atlas/symbol_game/tools/replay_formal_math_baselines.py}
  and its record
  \nolinkurl{atlas/coda_atlas/symbol_game/verification/formal_math_baselines.json}.
  It reads the runners and checkers from commits
  \nolinkurl{ea2b55d841bdbe283f2f6ac77abe8a1677cbf0f8} (the C045 runner)
  and \nolinkurl{46bd1dec156260f073b77caa07d8a453f8efad30} (C046 to
  C050). The C045 checkers are read from commit
  \nolinkurl{c568c20062697b54ddb8729e2476d4cfdb5c4982}, which differs
  from the run commit in that folder only by one import line that Python
  3.12 needs to load the exact checker.
\item
  Scoring receipts:
  \nolinkurl{c043_ramsey_vdw_constructions/commit_reveal/c043_l41_scoring_receipt.json}
  (commit \nolinkurl{50197478726e5ecb791e519760cee9a79af2af27}) and
  \nolinkurl{c044_finite_group_identification/commit_reveal/c044_l41_scoring_receipt.json}
  (commit \nolinkurl{999b36c1afed0f905f634f83adc1c646da7b975e}).
\end{itemize}

\subsection{5. Limitations}\label{limitations}

The C017 pilot has five cases, one reply per case and arm, two models,
and arms that change several things at once. The formal-mathematics
pilots have ten cases per cell and one reply per case, and their prompts
named the expected result for 36 of the 60 cases whose prompts can be
reconstructed (Table 4). Neither pilot supports a claim that one prompt
or model is better than another, or any claim about unaided capability,
and none is made. Replies from hosted models can vary between calls even
at temperature 0.

One of the three models is undisclosed and may not remain available. Two
of the run reports (C017 Arm 5 with
\nolinkurl{meta-llama/llama-3.3-70b-instruct} and C050 run 1) are
retained only in the repository history, and the C043 and C044 run
reports and the C049 panel B runner are missing. The two repeated runs
in Table 3 show that a single reply per case can change between runs.

The package replays only the C017 cases. The C045 to C050 checkers are
replayed by a script of the collection, outside this package, which
needs a repository checkout with the cited commits; the observations
about the models' answers still rest on the retained checker messages.
The C017 harness checkers share code with the challenge's reference
solver. The Lean file formalizes only the remark of Section 2.3; the
flatness check itself is tested, not mechanically formalized.

This paper assigns no authorization and changes no jurisdiction. It
carries no operational authority or institutional standing.

\subsection{6. Reproduction}\label{reproduction}

From the package root, with Python 3.12.10, pytest 8.3.5 and ruff
0.14.13:

\begin{lstlisting}
py -3.12 -m pytest code/tests/test_holonomy_checker.py -q
py -3.12 -m pytest code/tests/test_c017_replay.py -q
py -3.12 code/replay_c017.py
py -3.12 -m ruff check code
\end{lstlisting}

The Lean lemma builds with \passthrough{\lstinline!lake build!} inside
\passthrough{\lstinline!lean4/!}, using the toolchain pinned in
\passthrough{\lstinline!lean4/lean-toolchain!} (Lean 4.34.0-rc1);
\nolinkurl{tools/replay_lean_axiom_census.py} replays the axiom report.
The tests read only files inside the package and make no network call.
The model runs are not repeated here; their records are listed in
Sections 3.3 and 4.4. The replay behind Table 4 runs from the collection
folder of a repository checkout,
\passthrough{\lstinline!py -3.12 tools/replay\_formal\_math\_baselines.py --check!},
and compares its output with the recorded file.

\subsection{Verification for this
revision}\label{verification-for-this-revision}

The package verifier (\passthrough{\lstinline!verify\_deposit.py!})
checks the manifest and the two-way file census, and
\passthrough{\lstinline!tools/rebuild\_release.py --check!} checks the
release archive against the manifests. The package tests, lint, and Lean
build are executed by the dated review receipt recorded with this
revision; that receipt lists each command, tool version, exit status,
and tested scope. Historical receipts in this package keep their
original dates and scope and are not current execution evidence.

\subsection{References}\label{references}

\begin{enumerate}
\def\labelenumi{\arabic{enumi}.}
\tightlist
\item
  K. G. Wilson, ``Confinement of quarks,'' \emph{Physical Review D}
  \textbf{10}(8), 2445--2459 (1974).
  \href{https://doi.org/10.1103/PhysRevD.10.2445}{\nolinkurl{doi:10.1103/PhysRevD.10.2445}}
\item
  J. B. Kogut, ``An introduction to lattice gauge theory and spin
  systems,'' \emph{Reviews of Modern Physics} \textbf{51}(4), 659--713
  (1979).
  \href{https://doi.org/10.1103/RevModPhys.51.659}{\nolinkurl{doi:10.1103/RevModPhys.51.659}}
\item
  A. Hatcher, \emph{Algebraic Topology}. Cambridge University Press,
  2002.
\item
  D. E. Knuth and P. B. Bendix, ``Simple word problems in universal
  algebras,'' in J. Leech (ed.), \emph{Computational Problems in
  Abstract Algebra}, pp.~263--297. Pergamon Press, 1970.
  \href{https://doi.org/10.1016/B978-0-08-012975-4.50028-X}{\nolinkurl{doi:10.1016/B978-0-08-012975-4.50028-X}}
\item
  E. H. Bareiss, ``Sylvester's identity and multistep integer-preserving
  Gaussian elimination,'' \emph{Mathematics of Computation}
  \textbf{22}(103), 565--578 (1968).
  \href{https://doi.org/10.1090/S0025-5718-1968-0226829-0}{\nolinkurl{doi:10.1090/S0025-5718-1968-0226829-0}}
\end{enumerate}

\end{document}
